\section{Discussion}

The expanded manuscript supports a narrower and stronger claim than a generic ``agents beat RAG'' story. Classical RAG remains a credible default when the educational task is primarily evidence-grounded and success depends on retrieving the right passages efficiently. The repository-grounded snapshots nevertheless suggest that tutoring quality often depends on a second axis that classical RAG does not model directly: pedagogical coordination. On EduBench, the strongest configuration in the aligned snapshot is the non-RAG multi-agent route, which improves rubric coverage and score alignment while also reducing latency relative to classical RAG. On TutorEval, the strongest configuration is hybrid-fast agentic RAG, which preserves retrieval and improves both overlap-style metrics and response speed relative to the classical baseline. The preferred architecture therefore shifts with benchmark demands rather than remaining fixed.

That shift has practical implications for educational system design. If a deployment target mainly requires source-grounded answers over course materials, a well-tuned classical RAG stack may still be the simplest and strongest choice. If the target requires misconception diagnosis, scaffold calibration, rubric-sensitive critique, or adaptive next-step planning, the main system problem becomes coordination among instructional roles rather than retrieval alone. In those cases, a planner, diagnoser, rubric agent, tutor, and critic may earn their complexity. The architecture question is therefore not whether retrieval disappears, but whether retrieval remains the top-level controller.

This framing also clarifies why hybrid designs are plausible endpoints. The TutorEval snapshot suggests that retrieval is still valuable when long-context evidence matters, yet the best route is not fixed classical RAG. Instead, retrieval appears strongest when embedded inside a broader policy that decides when evidence is necessary and how it should constrain tutoring behavior. That is a more disciplined conclusion than claiming that non-RAG systems dominate. It also aligns with the broader educational intuition that a useful tutor must know both when to ground in sources and when to shift into diagnosis, questioning, or scaffolded explanation.

Several limitations remain important. First, the repository results are aligned early-run snapshots rather than full-benchmark evaluations, so they support directional evidence rather than definitive leaderboard claims. Second, the benchmarks expose different supervision signals: EduBench emphasizes rubric-style pedagogical scoring, while TutorEval more directly rewards evidence-grounded tutoring behavior. Third, some metrics are architecture-sensitive by construction, especially grounded overlap in runs that do not retrieve. These are reasons to narrow the paper's claims, not to discard the comparison. They explain why the article positions the current evidence as repository-grounded support for a broader evaluation framework rather than as final proof of classroom superiority.

The paper also needs to remain careful about what benchmark wins mean. Even strong performance on educational benchmarks is not the same as evidence of improved learner outcomes. Human evaluation can assess scaffold quality, tone, next-step usefulness, and level appropriateness, but not long-term learning gains. The manuscript therefore keeps pedagogical quality, groundedness, adaptivity, and cost separate and treats learner-impact claims as future work requiring longitudinal study.

Finally, negative or mixed results are theoretically useful. If full evaluations eventually show that classical RAG remains dominant on simple educational tasks, that would validate the taxonomy's low-coordination regime. If retrieval-enabled hybrids outperform non-RAG systems on fact-intensive tutoring, that would clarify where grounding is non-negotiable. The point of the framework is not to guarantee that more coordination wins. It is to make architecture choice legible, benchmark by benchmark, so that educational AI systems are compared according to the kind of difficulty they actually face.
